\SourceLeaf{61}
\begin{center}
\textit{FÜNFTER BEWEIS DES FUNDAMENTALTHEOREMS}\par
\textit{IN DER THEORIE DER QUADRATISCHEN RESTE.}
\end{center}

\Article{1}

In der Einleitung haben wir bereits gesagt, dass der fünfte und der dritte Beweis von demselben Lemma ausgehen. Der Bequemlichkeit halber erschien es uns angebracht, es hier in einer für die vorliegende Untersuchung passenden Bezeichnung zu wiederholen.

\textsc{Lemma.}\quad \textit{Es sei $m$ eine Primzahl (positiv und ungerade) und $M$ eine ganze, durch $m$ nicht teilbare Zahl. Man nehme die kleinsten positiven Reste der Zahlen}

\[
\begin{aligned}
M,\ 2M,\ 3M,\ 4M\ \mathinner{.\,.\,.\,.\,.\,.}\ \tfrac12(m-1)M
\end{aligned}
\]

\textit{modulo $m$. Diese sind teils kleiner als $\tfrac12m$, teils größer; die Anzahl der Letzteren sei $=n$. Dann ist $M$ ein quadratischer Rest oder Nichtrest von $m$, je nachdem, ob $n$ gerade oder ungerade ist.}

{\rightskip=0pt plus1fil \textsc{Beweis.}\quad Diejenigen unter diesen Resten, die kleiner als $\tfrac12m$ sind, seien $a,b,c,d$ usw., die übrigen, die größer als $\tfrac12m$ sind, dagegen $a',\ b',\ c',\ d'$ usw. Die Ergänzungen der Letzteren zu $m$, nämlich \mbox{$m-a',\ m-b',\ m-c',\ m-d'$} usw., sind offenbar sämtlich kleiner als $\tfrac12m$ und sowohl untereinander als auch von den Resten $a,\ b,\ c,\ d$ usw. verschieden. Zusammen mit diesen bilden sie daher, wenn auch in anderer Reihenfolge, genau alle Zahlen $1,\ 2,\ 3,\ 4\dotsfour\tfrac12(m-1)$. Setzen wir also das Produkt\par}

\[
\begin{aligned}
1\pdot2\pdot3\pdot4\dotsfive\tfrac12(m-1)=P
\end{aligned}
\]

so gilt

\[
\begin{aligned}
P=abcd\dotsfour\times(m-a')(m-b')(m-c')(m-d')\dotsfour
\end{aligned}
\]

und damit

\[
\begin{aligned}
(-1)^nP=abcd\dotsfour\times(a'-m)(b'-m)(c'-m)(d'-m)\dotsfour
\end{aligned}
\]

Ferner gilt modulo $m$

\[
\begin{aligned}
PM^{\frac12(m-1)}&\equiv abcd\dotstwo\times a'b'c'd'\dotstwo\\&\equiv abcd\dotstwo\times(a'-m)(b'-m)(c'-m)(d'-m)\dotstwo
\end{aligned}
\]

und damit

\[
\begin{aligned}
PM^{\frac12(m-1)}\equiv P(-1)^n
\end{aligned}
\]

Hieraus folgt $M^{\frac12(m-1)}\equiv\pm1$, wobei das obere oder untere Vorzeichen zu nehmen ist, je nachdem, ob $n$ gerade oder ungerade ist. Mithilfe des in den \textit{Disquisitiones Arithmeticae}, Art. 106, bewiesenen Satzes ergibt sich daraus unmittelbar die Richtigkeit des Lemmas.

\SourceLeaf{62}
\Article{2}

\textsc{Satz.}\quad \textit{Es seien $m,M$ positive ungerade ganze Zahlen, die zueinander teilerfremd sind, und $n$ sei die Anzahl derjenigen unter den kleinsten positiven Resten der Zahlen}

\[
\begin{aligned}
M,\ 2M,\ 3M\dotsfive\tfrac12(m-1)M
\end{aligned}
\]

\textit{modulo $m$, die größer als $\frac12m$ sind. Ebenso sei $N$ die Anzahl derjenigen unter den kleinsten positiven Resten der Zahlen}

\[
\begin{aligned}
m,\ 2m,\ 3m\dotsfive\tfrac12(M-1)m
\end{aligned}
\]

\textit{modulo $M$, die größer als $\frac12M$ sind.\quad Dann sind die drei Zahlen $n,N$, $\frac14(m-1)(M-1)$ entweder sämtlich gerade, oder eine ist gerade und die beiden anderen sind ungerade.}

\textsc{Beweis.}\quad Wir bezeichnen

{\small
\begin{center}\setlength{\tabcolsep}{3pt}
\begin{tabular}{@{}lll@{}}
mit $f$ & die Sammlung der Zahlen & $1,\ 2,\ 3\dotsfive\frac12(m-1)\mathord{.}$ \\
mit $f'$ & die Sammlung der Zahlen & $m-1,\ m-2,\ m-3\dotsfive\frac12(m+1)$ \\
mit $F$ & die Sammlung der Zahlen & $1,\ 2,\ 3\dotsfive\frac12(M-1)$ \\
mit $F'$ & die Sammlung der Zahlen & $M-1,\ M-2,\ M-3\dotsfive\frac12(M+1)$
\end{tabular}\end{center}
}

Damit gibt $n$ an, wie viele Zahlen $Mf$ ihre kleinsten positiven Reste modulo $m$ in der Sammlung $f'$ haben. Ebenso gibt $N$ an, wie viele Zahlen $mF$ ihre kleinsten positiven Reste modulo $M$ in der Sammlung $F'$ haben. Schließlich bezeichne

{\small
\begin{center}\setlength{\tabcolsep}{3pt}
\begin{tabular}{@{}lll@{}}
$\varphi$ & die Sammlung der Zahlen & $1,\ 2,\ 3\dotsfive\frac12(mM-1)$ \\
$\varphi'$ & die Sammlung der Zahlen & $mM-1,\ mM-2,\ mM-3\dotsfive\frac12(mM+1)$
\end{tabular}\end{center}
}

Jede ganze, durch $m$ nicht teilbare Zahl muss modulo $m$ entweder einem Rest aus $f$ oder einem Rest aus $f'$ kongruent sein. Ebenso ist jede ganze, durch $M$ nicht teilbare Zahl modulo $M$ entweder einem Rest aus $F$ oder einem Rest aus $F'$ kongruent. Daher lassen sich alle Zahlen in $\varphi$, unter denen offenbar keine zugleich durch $m$ und $M$ teilbar ist, folgendermaßen in acht Klassen einteilen.

I.\quad Die erste Klasse enthält die Zahlen, die modulo $m$ einer Zahl aus $f$, modulo $M$ aber einer Zahl aus $F$ kongruent sind.\quad Ihre Anzahl bezeichnen wir mit $\alpha$.

\SourceLeaf{63}
II.\quad Die Zahlen, die modulo $m,M$ jeweils Zahlen aus $f,F'$ kongruent sind; ihre Anzahl setzen wir $=\beta$.

III.\quad Die Zahlen, die modulo $m,M$ jeweils Zahlen aus $f',F$ kongruent sind; ihre Anzahl setzen wir $=\gamma$.

IV.\quad Die Zahlen, die modulo $m,M$ jeweils Zahlen aus $f',F'$ kongruent sind; ihre Anzahl sei $=\delta$.

V.\quad Die Zahlen, die durch $m$ teilbar, modulo $M$ aber Resten aus $F$ kongruent sind.

VI.\quad Die Zahlen, die durch $m$ teilbar, modulo $M$ aber Resten aus $F'$ kongruent sind.

VII.\quad Die Zahlen, die durch $M$ teilbar, modulo $m$ aber Resten aus $f$ kongruent sind.

VIII.\quad Die Zahlen, die durch $M$ teilbar, modulo $m$ aber Resten aus $f'$ kongruent sind.

Offenbar umfassen die Klassen V und VI zusammen alle Zahlen $mF$. Die Anzahl der in VI enthaltenen Zahlen ist $=N$, also ist die Anzahl der in V enthaltenen Zahlen $\frac12(M-1)-N$.\quad Ebenso enthalten die Klassen VII und VIII zusammen alle Zahlen $Mf$; in der Klasse VIII finden sich $n$ Zahlen, in der Klasse VII dagegen $\frac12(m-1)-n$.

Auf genau dieselbe Weise werden alle Zahlen in $\varphi'$ in acht Klassen IX\,.\,.\,XVI eingeteilt. Behalten wir dabei dieselbe Reihenfolge bei, so sieht man leicht, dass die Zahlen in den Klassen

\begin{center}IX, X, XI, XII, XIII, XIV, XV, XVI\end{center}

jeweils die Ergänzungen zu $mM$ der Zahlen in den Klassen

\begin{center}IV, III, II, I, VI, V, VIII, VII\end{center}

sind, so dass sich in der Klasse IX $\delta$ Zahlen, in der Klasse X $\gamma$ Zahlen usw. finden.\quad Nehmen wir nun alle Zahlen der ersten Klasse mit allen Zahlen der neunten zusammen, so erhalten wir offenbar alle Zahlen unterhalb $mM$, die modulo $m$ einer Zahl aus $f$, modulo $M$ aber einer Zahl aus $F$ kongruent sind. Deren Anzahl ist, wie man leicht sieht, gleich der Anzahl aller Kombinationen jeder einzelnen Zahl aus $f$ mit jeder einzelnen aus $F$.\quad Wir haben daher

\[
\begin{aligned}
\alpha+\delta=\tfrac14(m-1)(M-1)
\end{aligned}
\]

\SourceLeaf{64}
und aus ähnlichem Grund gilt auch

\[
\begin{aligned}
\beta+\gamma=\tfrac14(m-1)(M-1)
\end{aligned}
\]

\Needspace{9\baselineskip}
Nehmen wir alle Zahlen der Klassen II, IV und VI zusammen, so erhalten wir offenbar alle Zahlen unterhalb $\frac12mM$, die einem Rest aus $F'$ modulo $M$ kongruent sind.\quad Dieselben Zahlen lassen sich aber auch so darstellen:

\[
\begin{aligned}
F',\ M+F',\ 2M+F',\ 3M+F'\dotsfour\tfrac12(m-3)M+F'
\end{aligned}
\]

Ihre Gesamtzahl ist daher $=\frac14(m-1)(M-1)$, oder es gilt

\[
\begin{aligned}
\beta+\delta+N=\tfrac14(m-1)(M-1)
\end{aligned}
\]

Ebenso lässt sich aus der Vereinigung aller Zahlen der Klassen III, IV und VIII folgern:

\[
\begin{aligned}
\gamma+\delta+n=\tfrac14(m-1)(M-1)
\end{aligned}
\]

Aus diesen vier Gleichungen ergeben sich die folgenden:

\[
\begin{aligned}
2\alpha&=\tfrac14(m-1)(M-1)+n+N\\
2\beta&=\tfrac14(m-1)(M-1)+n-N\\
2\gamma&=\tfrac14(m-1)(M-1)-n+N\\
2\delta&=\tfrac14(m-1)(M-1)-n-N
\end{aligned}
\]

von denen jede die Richtigkeit des Satzes zeigt.

\Article{3}

Nehmen wir nun an, $m$ und $M$ seien Primzahlen, so ergibt sich das Fundamentaltheorem unmittelbar aus der Verbindung des vorhergehenden Satzes mit dem Lemma aus Art. 1.\quad Denn offenbar gilt:

I.\quad Wenn die Zahlen $m,M$ beide die Form $4k+1$ haben oder nur eine von ihnen, ist die Zahl $\frac14(m-1)(M-1)$ gerade. Daher sind $n$ und $N$ entweder beide gerade oder beide ungerade, also ist entweder jede der Zahlen $m$ und $M$ ein quadratischer Rest der anderen oder jede ein quadratischer Nichtrest der anderen.

II.\quad Wenn dagegen beide Zahlen $m,M$ die Form $4k+3$ haben, ist \mbox{$\frac14(m-1)(M-1)$} ungerade. Daher ist eine der Zahlen $n,N$ gerade und die andere ungerade; folglich ist eine der Zahlen $m,M$ ein quadratischer Rest der anderen und diese ein quadratischer Nichtrest der ersteren. Q. E. D.
